\documentclass[11pt]{article}
\usepackage[margin=1in]{geometry}
\usepackage{amsmath,amssymb,amsthm,hyperref}
\newtheorem{theorem}{Theorem}
\title{A triangularizing representation basis needs three vectors}
\author{ziangni-sys}
\date{Conjecture 00000008587}
\begin{document}
\sloppy
\maketitle
\noindent The source defines adaptive bases as minimal bases making the
representation matrices upper triangular and asserts that the dimension
upper bound of such minimal bases is $2$. We disprove this clause for a
commutative operator algebra in a concrete faithful unital representation.

\section*{The operator algebra and representation}
Let $H=\mathbb C^3$, equipped with its usual Euclidean Hilbert-space structure,
and let
\[
 \mathcal A=\{a I_H:a\in\mathbb C\}\subseteq\operatorname{End}_{\mathbb C}(H).
\]
This is a unital commutative complex operator algebra. The scalar map
$\rho:\mathbb C\to\operatorname{End}_{\mathbb C}(H)$, $a\mapsto a I_H$,
is an algebra homomorphism, since
\[
 \rho(a+b)=\rho(a)+\rho(b),\qquad
 \rho(ab)=\rho(a)\rho(b),\qquad \rho(1)=I_H.
\]
It is faithful: if $a I_H=b I_H$, evaluating at the first standard basis
vector gives $a=b$. Every such finite-dimensional linear operator is bounded.

For every vector-space basis $B$ of $H$, the matrix of $\rho(a)$ is $a I_3$.
Indeed, $\rho(a)b_j=ab_j$ for each basis vector $b_j$. The matrix is diagonal,
and therefore upper triangular. In particular the standard three-vector
basis is an actual triangularizing basis for every element of $\mathcal A$.

\begin{theorem}
The minimum cardinality of a triangularizing representation basis for
$\mathcal A$ acting on $H$ is $3$, contradicting the asserted upper bound $2$.
\end{theorem}
\begin{proof}
The standard basis provides a triangularizing basis of cardinality $3$.
Every vector-space basis of $H$ is finite and has cardinality
$\dim_{\mathbb C}H=3$. This applies to every possible triangularizing basis,
not just the standard one. Thus the minimum is exactly $3$, and no
triangularizing basis has cardinality at most $2$.
\end{proof}

\section*{Meaning of minimality and scope}
The bases here are vector bases of the space carrying the representation,
as required to make its matrices upper triangular. Minimality is with
respect to such representation bases for this fixed operator algebra
acting on this fixed space. No vector can be omitted while retaining a
basis of $H$.

We do not replace this notion by a basis of the abstract algebra
$\mathcal A$ itself, whose algebraic dimension is $1$, or minimize the
dimension over different representation spaces. The source defines bases
through representation matrices, and specifies no such alternative
minimization. Our witness uses an actual faithful unital representation,
and computes the actual representation dimension and all basis sizes.
The failed upper-bound clause is a conjunct; the other claims about maximal
ideal spaces and noncommutative algebras need not be resolved.

\newpage
\section*{Formal objects and correspondence}
The Lean space \texttt{H} is Mathlib's actual
\texttt{EuclideanSpace} $\mathbb C$ \texttt{(Fin 3)}, and \texttt{EndH}
is its complex linear endomorphism algebra.
The actual algebra homomorphism \texttt{scalar} is
\texttt{Algebra.ofId} into this endomorphism algebra.
\texttt{scalar\_apply} verifies $\rho(a)x=ax$,
\texttt{scalar\_injective} verifies faithfulness,
\texttt{scalar\_unital} verifies the unit identity, and
\texttt{scalar\_continuous} verifies continuity of every represented scalar.
\texttt{scalarAlgebra} is the actual range subalgebra.
\texttt{algebra\_commutative} proves that every pair of its elements commutes.

\texttt{standardBasis} is the actual standard orthonormal basis viewed as
a vector-space \texttt{Basis}.
\texttt{scalar\_matrix} uses Mathlib's actual
\texttt{LinearMap.toMatrix} and proves that the representation matrix is
$aI$ in every genuine finitely indexed basis.
\texttt{Triangularizing} quantifies over every algebra element and requires
each entry strictly below the diagonal to be zero.
\texttt{every\_basis\_triangularizes} proves this property in any such
ordered basis, and \texttt{standard\_triangularizes} supplies its explicit
standard-basis instance.

\texttt{ambient\_dimension} computes the actual complex
\texttt{Module.finrank} as $3$. \texttt{every\_basis\_finite} proves that
every basis index type is finite, rather than assuming this in the
minimality argument. \texttt{every\_basis\_cardinality} proves that the
actual \texttt{Nat.card} of every basis index type is $3$, including
arbitrary index universes. Finite-dimensionality makes this cardinal
computation an ordinary finite basis count.

The theorem \texttt{minimum\_triangularizing\_basis\_size} combines an
actual triangularizing basis with the universal cardinality equality.
\texttt{dimension\_bound\_fails} combines the actual standard basis with
the impossibility of any basis having cardinality at most $2$.
No representation dimension, basis size or matrix data is assumed.

\section*{Reproduction and validation}
The project pins Lean 4.19.0 and Mathlib commit
\begin{center}\small\texttt{c44e0c8ee63ca166450922a373c7409c5d26b00b}.\end{center}
From \texttt{lean/}, run \texttt{lake update}, \texttt{lake exe cache get},
and \texttt{lake build}; check the entire source with
\begin{center}\small\texttt{lake env lean -DwarningAsError=true Main.lean}.\end{center}
The source prints dependency audits for the actual representation,
commutativity, matrices, triangularization, basis finiteness, cardinality
and final contradiction. Only standard logical axioms are used. There are
no admitted statements, custom axioms, native decision procedures or
auxiliary numerical evidence.
\end{document}
